\subsubsection*{Solution}
Let

\[
a:=\cos ^2\theta,\qquad b:=\sin ^2\theta=1-a,
\]

and use base-$2$ logarithms, so the Holevo information is measured in bits. Define the binary entropy

\[
h_2(t):=-t\log_2t-(1-t)\log_2(1-t),
\]

with $0\log_2 0:=0$. The Holevo quantity of an ensemble is

\[
\chi=S\left(\sum_jp_j\rho_j\right)-\sum_jp_jS(\rho_j),
\]

as given in standard quantum-information references such as \href{\detokenize{https://www.preskill.caltech.edu/ph229/notes/chap5.pdf}}{Preskill's lecture notes}. The von Neumann entropy is computed from the eigenvalues as $S(\rho)=-\operatorname{Tr}(\rho\log_2\rho)$; see also the \href{\detokenize{https://quantum.cloud.ibm.com/docs/api/qiskit/0.45/quantum_info}}{IBM Quantum documentation}.

{
\color{blue}
The AI agent introduces the objects that need to be computed to solve the problem. Interestingly, it defines the Holevo information using the greek variable $\chi$ as on the Wikipedia page (and I assume in Preskill's notes). However I would consider this a \textbf{typo} as $\chi$ was used in the problem statement to define the initial classical-quantum state and not the Holevo information. 

}

\paragraph*{Step-by-Step Derivation}

\subparagraph*{1. Entropy of each individual state}

The lower-right $2\times2$ block of $\rho_j$ is

\[
B_j=
\begin{pmatrix}
(1-\gamma_j)a &
\sqrt{1-\gamma_j}\cos\theta\sin\theta\,e^{i\phi_j}\\
\sqrt{1-\gamma_j}\cos\theta\sin\theta\,e^{-i\phi_j}&b
\end{pmatrix}.
\]

Its determinant is

\[
\det B_j =(1-\gamma_j)ab -(1-\gamma_j)\cos^2\theta\sin^2\theta =0,
\]

while

\[
\operatorname{Tr}B_j=(1-\gamma_j)a+b=1-a\gamma_j.
\]

Therefore, the eigenvalues of $B_j$ are $0$ and $1-a\gamma_j$. Including the upper-left eigenvalue $a\gamma_j$, the spectrum of $\rho_j$ is

\[
\operatorname{spec}(\rho_j) = \left\{a\gamma_j,\ 1-a\gamma_j,\ 0\right\}.
\]

Hence

\[
S(\rho_j)=h_2(a\gamma_j).
\]

{
\color{blue}
The AI agent finds the eigenvalues correctly (it is less lazy than me in that regard).
}

\subparagraph*{2. Form of the ensemble-average state}

Introduce

\[
\bar\gamma:=\sum_jp_j\gamma_j, \qquad u:=\sum_jp_j\sqrt{1-\gamma_j}\,e^{i\phi_j}.
\]

Then

\[
\bar\rho:=\sum_jp_j\rho_j
=
\begin{pmatrix}
a\bar\gamma&0&0\\
0&a(1-\bar\gamma)&\cos\theta\sin\theta\,u\\
0&\cos\theta\sin\theta\,u^*&b
\end{pmatrix}.
\]

Thus all phase dependence occurs through the single complex number $u$.

{
\color{blue}
This step in the derivation is not well motivated, so it might confuse the reader onto why one need this. However, one will see that $\bar\gamma$ is the crucial part for the solution. 
}

\subparagraph*{3. Removing the average coherence increases entropy}

Let

\[
D=\operatorname{diag}(1,1,-1).
\]

The map

\[
\Delta(\bar\rho) = \frac{1}{2}\left(\bar\rho+D\bar\rho D^\dagger\right)
\]

removes the off-diagonal entries:

\[
\Delta(\bar\rho) = \operatorname{diag}\left( a\bar\gamma,\, a(1-\bar\gamma),\, b \right).
\]

By concavity and unitary invariance of the von Neumann entropy,

\[
S(\Delta(\bar\rho)) \geq \frac12S(\bar\rho)+\frac12S(D\bar\rho D^\dagger) =S(\bar\rho).
\]

These entropy properties, including the fact that deleting off-diagonal entries cannot decrease entropy, are summarized in \href{\detokenize{https://www.preskill.caltech.edu/ph229/notes/chap5.pdf}}{Preskill's notes}.

Consequently,

\[
S(\bar\rho) \leq H_3\!\left( a\bar\gamma,\, a(1-\bar\gamma),\, b \right),
\]

where

\[
H_3(q_1,q_2,q_3):=-\sum_{k=1}^3q_k\log_2q_k.
\]

{
\color{blue}
This is an interesting move from the AI agent: it decides to bound the entropy instead of treating it in full. As we will see later this is it's way for resolve the question's ambiguity (without telling us).
}

\subparagraph*{4. Minimize the average individual-state entropy}

Because $h_2$ is concave, the function

\[
g(\gamma):=h_2(a\gamma)
\]

is concave on $[0,1]$. A concave function lies above the chord joining its endpoint values, so

\[
h_2(a\gamma) \geq (1-\gamma)h_2(0)+\gamma h_2(a) = \gamma h_2(a).
\]

Averaging gives

\[
\sum_jp_jS(\rho_j) = \sum_jp_jh_2(a\gamma_j) \geq \bar\gamma\,h_2(a).
\]

Combining the two bounds,

\[
\chi \leq H_3\!\left( a\bar\gamma,\, a(1-\bar\gamma),\, b \right) -\bar\gamma h_2(a).
\]

{
\color{blue}
Here, it uses the bound on entropy to bound the Holevo information.
}


\subparagraph*{5. Reduce to one scalar variable}

Set

\[
x:=\bar\gamma\in[0,1].
\]

{
\color{blue}
This is the very crucial and "smart" step in the AI agent's derivation. It avoids the ambiguity by reducing the problem to only one variable that is $\bar\gamma$ which is the average the mean value of the random variable $\gamma_x$ sampled according to the probability distribution $\{p_x\}$. Neither of them is known but the AI agent decided to assume that $\bar\gamma$ can be computed and thus uses it instead. The reduction is possible thanks to dropping the coherences in step 3 when bounding the entropy. 
}


The ternary entropy simplifies as

\[
\begin{aligned}
H_3\left(ax,a(1-x),b\right)
&=-ax\log_2(ax)-a(1-x)\log_2(a(1-x))-b\log_2b\\
&=-a\log_2a-b\log_2b
-a\left[x\log_2x+(1-x)\log_2(1-x)\right]\\
&=h_2(a)+a\,h_2(x).
\end{aligned}
\]

Therefore,

\[
\chi\leq f(x),
\]

where

\[
f(x)=h_2(a)+a\,h_2(x)-x\,h_2(a).
\]

Restoring $a=\cos^2\theta$,

\[
{ f(x) = h_2(\cos^2\theta) +\cos^2\theta\,h_2(x) -x\,h_2(\cos^2\theta) }.
\]

\subparagraph*{6. Show that the bound is attainable}

For any fixed $x\in[0,1]$, take the three-state ensemble

\[
\begin{array}{c|c|c}
\text{Probability}&\gamma&\phi\\ \hline
x&1&0\\
\dfrac{1-x}{2}&0&0\\
\dfrac{1-x}{2}&0&\pi
\end{array}.
\]

The two $\gamma=0$ states have opposite coherences, so their off-diagonal entries cancel. Hence

\[
\bar\rho = \operatorname{diag}\left(ax,a(1-x),b\right).
\]

Moreover,

\[
\sum_jp_jS(\rho_j) = x\,h_2(a) +\frac{1-x}{2}h_2(0) +\frac{1-x}{2}h_2(0) = x\,h_2(a).
\]

Thus this ensemble achieves $\chi=f(x)$ for every $x$, proving that the upper bound is exact.

For completeness, when $a>0$, $f$ is strictly concave and its maximizing point is

\[
x_*= \frac{1}{1+2^{h_2(a)/a}}.
\]

If $a=0$, then $f(x)=0$ for all $x$.

{
\color{blue}
Here, it finds a state for which the bound is tight as a sanity check. Though, two things bother me: 
\begin{enumerate}
  \item I believe here it confused itself with what the variable $x$ is, since in the table and the two following equations it refers to the probability $p_x$ in defining the classical state (which is now a third definition of $x$). However, in the previous part it defined $x=\bar\gamma$ which are both quite different objects. I consider this to be a \textbf{typo} or \textbf{inconsistency}.
  
  \item The part that follows the "For completeness, ..." is in general correct but again this does not coincide with the definition of $x$ from above and no justification is provided. So I consider this to be an \textbf{inconsistency}.
\end{enumerate} 
}